\section{Introduction}
\label{sec:introduction}
AI systems can fail differently across tasks even when the same model is used. In our single-turn adversarial prompt campaign, financial-domain hallucination and in-context injection are elicited in $85.1\%$ and $43.1\%$ of classifiable responses, while jailbreak and harmful-content probes elicit at $0.7\%$ and $3.2\%$. The campaign does not test vendor safety scores or deployed agents. It motivates retaining category-level evidence rather than collapsing this particular suite into a single average.

Insurance provides one possible use for such evidence: underwriting can compare documented controls, experience adjustment can incorporate verified incident records, and claims auditing can check whether disclosed records match earlier commitments~\cite{romanosky_2019_content,eling_2016_cyberrisk}. None of these uses follows directly from prompt elicitation rates. Pricing additionally requires exposure denominators, loss severity, policy coverage, dependence across insureds, and calibration to actual deployments.

We examine three components of that evidence workflow: measurement, private pooling, and attestation. The components reuse existing adversarial-evaluation methods, generalization and diversity constraints, and Ed25519--Merkle commitments. Their contribution is a concrete implementation and a set of component evaluations, not new cryptographic primitives or an actuarial pricing model.

\subsection{Contributions}
\begin{enumerate}
\item \textbf{Measurement.} Twelve prompt categories, five hosted models, twenty prompts per category and three repeats produce $3{,}600$ trials. We report classifiable-response elicitation rates, missingness, and inter-judge agreement, distinguishing these from deployment incident rates.
\item \textbf{Pooling.} A synthetic-portfolio comparison measures suppression, ranking utility, and two specified attack simulators for suppression, hierarchical generalization, distinct $\ell$-diversity, and a Laplace-noise ablation. The results establish behavior in the tested generator and attack models; they do not establish anonymity against arbitrary auxiliary information.
\item \textbf{Attestation.} A Rust Ed25519--Merkle implementation and Python reference support batched integrity checks. We report measured signing throughput and $O(\log n)$ membership-proof verification costs, while identifying unsigned metadata and untested identity/time infrastructure.
\end{enumerate}

The public artifact is \url{https://github.com/sarkar-dipankar/agent-risk-evidence}. It includes recorded results, prompts, code, and an audit of manuscript-to-artifact consistency. No connected multi-turn agent or real insurance portfolio is evaluated.
